\documentclass[11pt]{article}
\usepackage{amsmath,amssymb,amsthm}

\newtheorem{theorem}{Theorem}
\newtheorem{corollary}[theorem]{Corollary}

\newcommand{\Up}{U_+}
\newcommand{\mass}{\operatorname{B}}
\newcommand{\car}{\operatorname{C}}
\newcommand{\sh}{\operatorname{D}}

\title{Exact carry-position identities for deep inverse prefixes}
\author{}
\date{27 August 2026}

\begin{document}
\raggedbottom
\maketitle

For a positive odd integer \(x\), put
\[
 \Up(x):=\frac{3x+1}{2^{v_2(3x+1)}}.
 \tag{1}
\]
For a positive exponent word \(E=(e_1,\ldots,e_h)\), define
\[
 \mass(E):=\sum_{i=1}^h e_i,
 \qquad
 \car(E):=\sum_{i=1}^h
 3^{h-i}2^{e_1+\cdots+e_{i-1}},
 \tag{2}
\]
and
\[
 \sh(E):=\car(E)+2^{\mass(E)}-3^h.
 \tag{3}
\]
An exact path from \(z\) to \(r\) with word \(E\) satisfies
\[
 2^{\mass(E)}(r+1)=3^h(z+1)+\sh(E).
 \tag{4}
\]

Write
\[
 e_i=1+c_i,
 \qquad
 c_i\geq0,
 \qquad
 t:=\sum_{i=1}^h c_i,
 \tag{5}
\]
and define the cumulative excess
\[
 p_k:=\sum_{i=1}^k c_i
 \qquad(0\leq k\leq h).
 \tag{6}
\]
Thus \(p_0=0\) and \(p_h=t\).

\section{The two exact slack identities}

For fixed \(h\) and \(t\), the shifted-carry extrema are
\[
 \sh_{\min}=2^h(2^t-1),
 \qquad
 \sh_{\max}=2\,3^{h-1}(2^t-1).
 \tag{7}
\]

\begin{theorem}[carry-position decomposition]
Every positive exponent word of length \(h\) and excess \(t\) satisfies
\[
 \boxed{
 \frac{\sh(E)-\sh_{\min}}{3^h}
 =
 \frac13\sum_{k=1}^{h-1}
 \left(\frac23\right)^k(2^{p_k}-1),}
 \tag{8}
\]
and
\[
 \boxed{
 \frac{\sh_{\max}-\sh(E)}{3^h}
 =
 \frac13\sum_{k=1}^{h-1}
 \left(\frac23\right)^k(2^t-2^{p_k}).}
 \tag{9}
\]
\end{theorem}

\begin{proof}
Equation (2) becomes
\[
 \car(E)=
 \sum_{i=1}^h3^{h-i}2^{i-1+p_{i-1}}.
 \tag{10}
\]
When all excess is in the final letter, every \(p_{i-1}\) in (10)
is zero.  Subtracting this minimum carry from (10), and putting
\(k=i-1\), gives
\[
 \sh(E)-\sh_{\min}
 =
 \sum_{k=1}^{h-1}
 3^{h-k-1}2^k(2^{p_k}-1).
 \tag{11}
\]
Division by \(3^h\) proves (8).

When all excess is in the first letter, every \(p_k\), \(k\geq1\), is
\(t\).  Subtracting (10) from this maximum carry gives
\[
 \sh_{\max}-\sh(E)
 =
 \sum_{k=1}^{h-1}
 3^{h-k-1}2^k(2^t-2^{p_k}).
 \tag{12}
\]
This proves (9).
\end{proof}

\section{The cumulative-excess squeeze}

Put
\[
 \mathcal A_h(r):=\left(\frac23\right)^h(r+1),
 \tag{13}
\]
and define the two sharp endpoints
\[
 L_{h,t}(r)
 :=2^t\left(\mathcal A_h(r)-\frac23\right)-\frac13,
 \tag{14}
\]
\[
 U_{h,t}(r)
 :=\left(\frac23\right)^h(2^t r+1)-1.
 \tag{15}
\]

\begin{corollary}[exact endpoint slack]
If \(E\) is an exact path from \(z\) to \(r\), then
\[
 \boxed{
 U_{h,t}(r)-z
 =
 \frac13\sum_{k=1}^{h-1}
 \left(\frac23\right)^k(2^{p_k}-1),}
 \tag{16}
\]
\[
 \boxed{
 z-L_{h,t}(r)
 =
 \frac13\sum_{k=1}^{h-1}
 \left(\frac23\right)^k(2^t-2^{p_k}).}
 \tag{17}
\]
\end{corollary}

\begin{proof}
Equation (4) gives
\[
 z+1=2^t\mathcal A_h(r)-\frac{\sh(E)}{3^h}.
 \]
Subtract the identities obtained by replacing \(\sh(E)\) with the two
extrema in (7), and apply Theorem~1.
\end{proof}

\begin{corollary}[two-sided position squeeze]
For every \(1\leq k<h\),
\[
 \boxed{
 \frac23\left(1-\left(\frac23\right)^k\right)
 (2^t-2^{p_k})
 \leq z-L_{h,t}(r),}
 \tag{18}
\]
\[
 \boxed{
 \left(\frac23\right)^k
 \left(1-\left(\frac23\right)^{h-k}\right)
 (2^{p_k}-1)
 \leq U_{h,t}(r)-z.}
 \tag{19}
\]
\end{corollary}

\begin{proof}
The sequence \(p_i\) is nondecreasing.  In (17), retain the terms
\(1\leq i\leq k\) and use
\[
 2^t-2^{p_i}\geq2^t-2^{p_k}.
 \]
The resulting geometric sum is (18).  In (16), retain the terms
\(k\leq i<h\) and use
\[
 2^{p_i}-1\geq2^{p_k}-1.
 \]
This gives (19).
\end{proof}

\begin{corollary}[finite screen for a bounded start]
Suppose that \(1\leq z<N\).  For every \(1\leq k<h\), put
\[
 \alpha_k:=\frac23\left(1-\left(\frac23\right)^k\right),
 \qquad
 \beta_k:=\left(\frac23\right)^k
 \left(1-\left(\frac23\right)^{h-k}\right).
 \tag{20}
\]
Then the integer \(p_k\) must satisfy
\[
 \boxed{
 2^{p_k}>
 2^t-\frac{N-L_{h,t}(r)}{\alpha_k},}
 \tag{21}
\]
\[
 \boxed{
 2^{p_k}\leq
 1+\frac{U_{h,t}(r)-1}{\beta_k}.}
 \tag{22}
\]
If these two bounds contain no power of two with exponent in
\(\{0,1,\ldots,t\}\), then no such prefix exists.
\end{corollary}

\begin{proof}
Use \(z<N\) in (18) and \(z\geq1\) in (19), and solve for
\(2^{p_k}\).  Since \(p_k\) is an integer between zero and \(t\), the
last assertion follows.
\end{proof}

The sharp all-excess interval uses only \(t\).  Equations (16)--(22)
retain where that excess occurs.  They give an exact finite obstruction
at every proposed triple \((h,t,r)\), but they do not yet exclude all
deep triples simultaneously.

\end{document}
